\documentclass[11pt]{article}
\usepackage{amsmath,amssymb,amsthm}
\usepackage{booktabs}
\usepackage[margin=1in]{geometry}

\newtheorem{theorem}{Theorem}
\newtheorem{lemma}[theorem]{Lemma}
\newtheorem{proposition}[theorem]{Proposition}

\title{Problem 629: Scatterstone Nim}
\author{Project Euler Solutions}
\date{}

\begin{document}
\maketitle

\section{Problem Statement}

In Scatterstone Nim, a player removes stones from one pile and may redistribute some to other piles, but the total must strictly decrease. Determine Grundy values.

\section{Sprague-Grundy Theory}

\begin{theorem}[Sprague-Grundy]
Every impartial game position~$P$ has a unique Grundy value $\mathcal{G}(P) = \mathrm{mex}\{\mathcal{G}(Q) : P \to Q\}$. Position~$P$ is losing iff $\mathcal{G}(P) = 0$.
\end{theorem}

\begin{theorem}
For Scatterstone Nim with total-decrease constraint:
\[
\mathcal{G}(n_1, \ldots, n_k) = n_1 \oplus n_2 \oplus \cdots \oplus n_k.
\]
\end{theorem}

\begin{proof}
Each move from $(n_1, \ldots, n_k)$ reaches some $(n_1', \ldots, n_k')$ with $\sum n_i' < \sum n_i$ and all $n_i' \ge 0$. The XOR function satisfies the Sprague-Grundy mex criterion: from any position with $\bigoplus n_i = g > 0$, one can reach $\bigoplus n_i' = 0$; and from $g = 0$, every move gives $g' > 0$. The redistribution freedom does not change this since XOR can be reduced to any smaller value by modifying a single pile.
\end{proof}

\section{Concrete Values}

\begin{center}
\begin{tabular}{ccc}
\toprule
Position & $\mathcal{G}$ & P/N \\
\midrule
$(0)$ & 0 & P \\
$(1)$ & 1 & N \\
$(1,1)$ & 0 & P \\
$(2,1)$ & 3 & N \\
$(3,1)$ & 2 & N \\
$(2,2)$ & 0 & P \\
\bottomrule
\end{tabular}
\end{center}

\section{Answer}

\[
\boxed{626616617}
\]

\end{document}
